\documentclass[journal]{IEEEtran}

%\usepackage[retainorgcmds]{IEEEtrantools}
%\usepackage{bibentry}  
\usepackage{xcolor,soul,framed}

\usepackage[pdftex]{graphicx}
\graphicspath{{../pdf/}{../jpeg/}}
\DeclareGraphicsExtensions{.pdf,.jpeg,.png}

\usepackage[cmex10]{amsmath}
\usepackage{amssymb}
\usepackage{amsfonts}
\newcommand{\mathbfcal}[1]{\boldsymbol{\mathcal{#1}}}

\usepackage{array}
\usepackage{mdwmath}
\usepackage{mdwtab}

\usepackage[numbers,sort&compress]{natbib}

\usepackage{eqparbox}
\usepackage{url}
\usepackage{caption}

% ---------------- ALGORITHM ----------------
\usepackage[ruled,vlined,linesnumbered]{algorithm2e}

% Define \Return used in your algorithm
\SetKw{Return}{return}

% -------------------------------------------

\usepackage{multirow}
\usepackage{booktabs}
\usepackage{float}
\usepackage{placeins}
\usepackage{dblfloatfix}
\usepackage{titlesec}
\usepackage{subcaption}

% Float-placement tuning for dense two-column layout.
% These settings keep figures/tables before the bibliography while reducing
% unnecessary white space caused by top-only float queues.
\setcounter{topnumber}{5}
\setcounter{bottomnumber}{5}
\setcounter{totalnumber}{10}
\setcounter{dbltopnumber}{5}
\renewcommand{\topfraction}{0.98}
\renewcommand{\bottomfraction}{0.98}
\renewcommand{\textfraction}{0.02}
\renewcommand{\floatpagefraction}{0.80}
\renewcommand{\dbltopfraction}{0.98}
\renewcommand{\dblfloatpagefraction}{0.80}
% Moderate float spacing for a compact two-column journal layout.
\setlength{\textfloatsep}{8pt plus 2pt minus 2pt}
\setlength{\dbltextfloatsep}{8pt plus 2pt minus 2pt}
\setlength{\floatsep}{7pt plus 2pt minus 2pt}
% Layout-only: use moderate, consistent spacing while preserving IEEEtran
% font, margins, and the mathematical content of all equations.
\setlength{\intextsep}{7pt plus 2pt minus 2pt}
\setlength{\dblfloatsep}{6pt plus 2pt minus 2pt}
\captionsetup[figure]{skip=3pt}
\captionsetup[table]{skip=3pt}
\AtBeginDocument{%
  \setlength{\abovedisplayskip}{4pt plus 1pt minus 1pt}%
  \setlength{\belowdisplayskip}{4pt plus 1pt minus 1pt}%
  \setlength{\abovedisplayshortskip}{2pt plus 1pt minus 1pt}%
  \setlength{\belowdisplayshortskip}{2.5pt plus 1pt minus 1pt}%
}



\begin{document}
\flushbottom
%\bstctlcite{IEEEexample:BSTcontrol}
   
\title{Multirate Physics-Guided Vehicle-Bridge Interaction Analysis with Spatiotemporal Traffic-Load Estimation and Occupancy-Load Consistency}


\author{\footnotesize
    Nidhi Agrawal\textsuperscript{1},
    Shashank Gupta\textsuperscript{1},
    Shuvendu Narayan Patel\textsuperscript{2},
    Devendra Gokul Patil\textsuperscript{3},
    Deeksha Arya\textsuperscript{4},
    Rajesh Kumar\textsuperscript{2}

    {\small
    \textsuperscript{1}Department of Computer Science and Information Systems, Birla Institute of Technology and Science, Pilani, Pilani Campus, India. \\
    \textsuperscript{2}Department of Civil Engineering, Birla Institute of Technology and Science, Pilani, Pilani Campus, India. \\
    \textsuperscript{3}Department of Mechanical Engineering, Birla Institute of Technology and Science, Pilani, Goa Campus, India. \\
    \textsuperscript{4}Center for Spatial Information Science, The University of Tokyo, Japan. \\
    }

    % Corresponding author
    {\scriptsize \textbf{Corresponding Author:} Shashank Gupta} \\[0.ex]
    \scriptsize
    \texttt{p20240429@pilani.bits-pilani.ac.in, shashank.gupta@pilani.bits-pilani.ac.in} \\
    \texttt{shuvendu@pilani.bits-pilani.ac.in, devendrap@goa.bits-pilani.ac.in} \\
    \texttt{deeksha@csis.u-tokyo.ac.jp, rajesh.kr@pilani.bits-pilani.ac.in}
}


\maketitle


\begin{abstract}
The complex dynamic interaction between stochastic traffic and long-span bridges has important implications for bridge design and safety assessment. Reliable vehicle--bridge interaction (VBI) analysis requires a numerically stable and consistent coupled formulation that captures both macroscopic traffic-flow conditions and microscopic vehicle dynamics. However, a crucial challenge lies in translating constantly changing traffic states into a physically consistent spatiotemporal load field for VBI analysis. Moreover, precise prediction of vehicle occupancy, spatial position, and speed does not necessarily ensure that the resulting traffic loads are spatially and temporally consistent with the underlying bridge domain. Such inconsistencies can propagate through the coupled formulation and result in errors in predicted bridge dynamic responses. This study proposes a physics-guided traffic--bridge interaction framework that utilizes the Intelligent Driver Model (IDM) to simulate longitudinal vehicle motion and a MOBIL-based lane-changing rule to capture microscopic traffic behaviour. The resulting physics-based traffic evolution is further refined through a learned spatiotemporal residual correction, followed by occupancy--load support enforcement to ensure spatial consistency between predicted vehicle occupancy and the corresponding traffic loads. The estimated traffic-load field is subsequently transferred to a reduced-order VBI model through multirate coupling to evaluate the resulting bridge dynamic response. Evaluation on 600 test windows from 37 unseen simulated traffic scenarios shows that enforcing occupancy--load support consistency reduces the maximum predicted midspan displacement from 191.12 to 40.61~mm and acceleration from 5.281 to 0.605~m/s$^2$. A frame-wise load-matched residual control further demonstrates that, despite essentially identical aggregate loading, the Support-Consistent Residual reduces mean midspan displacement RMSE from 4.395 to 3.727~mm and acceleration RMSE from 0.0439 to 0.0321~m/s$^2$, with positive 95\% scenario-cluster bootstrap confidence intervals for both improvements. These findings highlight that accurate traffic-state estimation alone is insufficient for reliable bridge-response prediction without maintaining spatial consistency between predicted vehicle occupancy and the corresponding traffic-load field.
\end{abstract}
\vspace{-4pt}
\begin{IEEEkeywords}
Vehicle-bridge interaction, traffic-load forecasting,
physics-guided learning, spatiotemporal residual correction,
reduced-order modeling, bridge dynamic response,
occupancy-load support consistency.
\end{IEEEkeywords}

\vspace{-3pt}
\section{Introduction}

Vehicle-bridge interaction (VBI) systems have been significantly investigated for characterizing the vibration response of vehicles and bridge structures under moving traffic load scenarios \cite{han2024,yin2025,huang2026,mei2025}. As traffic flow increases and vehicle loads become more diverse, the dynamic behaviour of the coupled system becomes increasingly complex \cite{aloisio2023,lifeng2023}. As a result, the loads imposed by moving vehicles and the resulting bridge response continuously influence each other through the VBI \cite{kreslin2024,shi2024,yau2024}. This complexity becomes more evident when multiple vehicles with distinct spatial positions, weights, velocities, and driving behaviours are simultaneously present on bridge structures \cite{yin2023,wei2025}. Consequently, realistic representation of traffic loading and accurate vehicle--bridge coupled analysis are essential for evaluating bridge dynamic response and structural safety under actual traffic conditions \cite{cantero2024,dong2023}.

To obtain precise bridge responses under stochastic traffic conditions, various traffic-flow models have been introduced to model the randomness and evolution of real-time traffic within vehicle-bridge coupled systems \cite{huang2026,wang2024}. The cellular automaton (CA) model has been widely utilized for traffic simulation since prescribed driving rules, including acceleration, deceleration, and lane-changing behaviour, can be represented in discrete space and time \cite{rickert1996,ruan2017}. However, the discrete-cell formulation imposes restrictions on vehicle spacing and velocity, which may limit the accurate modelling of stochastic traffic scenarios \cite{ruan2017}. In comparison with CA-based approaches, the Intelligent Driver Model (IDM) provides a deterministic, continuous space--time representation, which can describe individual vehicle positions and microscopic driving behaviour on the bridge \cite{treiber2000,geng2023,liu2023}. Existing IDM-based techniques have also incorporated lane-changing behaviour to simulate multi-lane traffic flow under different congestion scenarios \cite{caprani2016,obrien2015,xiong2020}.

Although IDM-based traffic simulation can reproduce realistic vehicle motion, the traffic state generated by the model does not, by itself, guarantee a physically consistent representation of the corresponding traffic-load field. In particular, the spatial distribution, temporal evolution, and loading support of individual vehicles must remain consistent with their predicted occupancy and motion to ensure reliable bridge-load estimation \cite{wangxu2019}. This traffic-to-load translation represents an important challenge in traffic--bridge interaction analysis. In previous vehicle bridge coupled studies, vehicle and bridge responses have often been solved as separate subsystems, with vehicle contact forces calculated in one environment and subsequently applied to the bridge model in another \cite{yangfonder1996,green1994,pan2002}. The bridge displacement response is then returned to the vehicle subsystem, and the interaction forces are recalculated iteratively until convergence is achieved. Recent traffic-bridge interaction frameworks have consequently sought to integrate traffic simulation information, vehicle models, and bridge models within a unified coupled system for dynamic response analysis.

In this context, Hamilton's principle has been increasingly adopted for the dynamic analysis of VBI systems since it provides an energy-based variational framework for the consistent formulation of the coupled vehicle-bridge dynamics \cite{mo2024}. In recent VBI formulations, the bridge has been represented using Euler-Bernoulli beam theory, while the vehicle subsystem is described using spring-damper models and coupled with the bridge through relative vehicle-bridge motion. The resulting coupled dynamic equations can subsequently be solved using the Newmark-$\beta$ method to obtain the time-dependent response of the VBI system \cite{mo2024,newmark1959}. 

The information obtained from the traffic-flow model can therefore be incorporated into the vehicle model to determine the traffic loading applied to the bridge. Compared with simplified moving-force or influence-line approaches, consideration of the vehicle-bridge coupled effect can provide a more realistic representation of bridge response under stochastic traffic flow \cite{yang2005,cantero2010}. Nevertheless, accurate estimation of traffic states alone may not be sufficient to ensure reliable bridge-response prediction if the corresponding load field does not preserve the required spatial and temporal consistency. This highlights the importance of explicitly addressing spatiotemporal traffic-load consistency within the traffic-bridge interaction framework.

Although the traffic state generated by the IDM can reproduce realistic vehicle motion, the initial spatial location, vehicle weight, and speed of all vehicles on the bridge are not provided by the IDM alone. More importantly, the influence of inconsistencies between predicted traffic occupancy and the corresponding spatial load distribution on subsequent bridge-response predictions remains insufficiently investigated. These limitations indicate the need for a unified framework that can incorporate realistic traffic information, preserve its spatiotemporal loading consistency, and subsequently propagate the resulting traffic loads into a coupled vehicle--bridge system for dynamic response analysis.

To address these challenges, this study develops a physics-guided traffic-bridge interaction framework for dynamic response analysis under evolving heterogeneous traffic scenarios. First, the traffic state is represented in terms of vehicle spatial location, velocity, weight, lane distribution, and driving behaviour, and its evolution is obtained using microscopic vehicle motion models. The Intelligent Driver Model (IDM) and the Minimizing Overall Braking Induced by Lane Changes (MOBIL) rule are employed to represent longitudinal vehicle motion and lane-changing behaviour, respectively. A residual spatiotemporal correction is then introduced to compensate for deviations in the physics-based traffic evolution and obtain the resulting traffic-load field, which is subsequently subjected to occupancy–load support consistency enforcement. In particular, the framework preserves consistency between predicted vehicle occupancy and the corresponding spatial support of bridge loading. Subsequently, the traffic-load information is incorporated into the vehicle model and coupled with the bridge model to establish the complete VBI system. Finally, the bridge and vehicle subsystems are represented using Euler-Bernoulli beam and spring--damper models, respectively, and the coupled dynamic response is obtained through Hamilton's principle and Newmark-$\beta$ time integration. The proposed framework is evaluated using unseen heterogeneous simulated traffic scenarios to investigate how occupancy–load support consistency within the spatiotemporal traffic-load field influences the predicted dynamic response of the bridge.


\begin{figure*}[!tbp]
    \centering
    \includegraphics[width=\textwidth]{Framework.png}
    \caption{Overall workflow of the proposed physics-guided traffic--bridge interaction framework.}
    \label{fig:overall_framew}
\end{figure*}

\section{Methodology}
\subsection{Overview of the Proposed Framework}
\label{sec:framework_overview}

The proposed framework addresses VBI under evolving heterogeneous traffic, where the structural response depends on the magnitude, spatial distribution, and temporal evolution of vehicle loading. Instead of treating traffic evolution and bridge dynamics independently, the framework interconnects the estimated traffic state directly to the subsequent VBI analysis through a spatially distributed traffic-load field.

Let the observed traffic history be denoted by $\mathbf{X}_{\mathrm{obs}}$, and let $\widehat{\mathbf{Y}}$ represent the estimated future traffic-load state. Hence, the overall response-estimation procedure can be expressed compactly as

\begin{equation}
\begin{aligned}
\widehat{\mathbf{R}}
&=
\mathcal{S}_{\mathrm{VBI}}
\left[
\mathcal{T}_{\mathrm{TB}}
\left(
\widehat{\mathbf{Y}}
\right)
\right],
\\
\widehat{\mathbf{Y}}
&=
\mathcal{F}_{\mathrm{tr}}
\left(
\mathbf{X}_{\mathrm{obs}}
\right).
\end{aligned}
\label{eq:overall_framework}
\end{equation}

where $\mathcal{F}_{\mathrm{tr}}$ denotes the traffic-state estimation process, $\mathcal{T}_{\mathrm{TB}}$ represents the traffic-to-bridge load transformation, and $\mathcal{S}_{\mathrm{VBI}}$ is the coupled VBI solution operator.

Unlike an aggregate-load representation, the estimated traffic field retains the lane-wise and longitudinal distribution of vehicle loading. The resultant bridge response is evaluated in terms of vertical displacement and acceleration at selected locations along the structural span. The complete workflow is illustrated in Fig.~\ref{fig:overall_framew}.

\subsection{Bridge-Centric Traffic State and Load-Field Representation}
\label{sec:traffic_state_representation}

\subsubsection{Vehicle-Level Traffic State}
\label{subsec:vehicle_state}
For vehicle $i$, the microscopic traffic state is
\begin{equation}
\mathbf{s}_i(t)=
[x_i(t),\,\ell_i(t),\,v_i(t),\,w_i,\,c_i]^{T},
\label{eq:vehicle_state}
\end{equation}
where $x_i$, $\ell_i$, $v_i$, $w_i$, and $c_i$ denote longitudinal position, lane, speed, vehicle weight, and vehicle category, respectively. The vehicles in the traffic domain form
$\mathcal{V}(t)=\{\mathbf{s}_i(t)\}_{i=1}^{N_v(t)}$, where $N_v(t)$ is the current vehicle count.

\subsubsection{Lane-Space Discretization}
\label{subsec:lane_space}
The roadway is divided into $N_\ell$ lanes and $N_x$ longitudinal cells of length $\Delta x=L_{\mathrm{tr}}/N_x$, where $L_{\mathrm{tr}}$ is the traffic-domain length. A vehicle is assigned to cell
\begin{equation}
s_i(t)=\min\!\left\{N_x,\,
1+\left\lfloor x_i(t)/\Delta x\right\rfloor\right\},
\label{eq:cell_index}
\end{equation}
and lane--cell location $\mathcal{C}_i(t)=(\ell_i(t),s_i(t))$. The implementation uses $N_\ell=L=4$ lanes and $N_x=S=50$ cells, preserving lane-wise and longitudinal positions for load transfer.

\subsubsection{Vehicle-State-to-Load-Field Mapping}
\label{subsec:load_field_mapping}
Let $\mathcal{V}_{\ell,s}(t)$ be the vehicles occupying lane $\ell$ and cell $s$. Occupancy count and binary support are
\begin{align}
O_{\ell,s}(t)&=|\mathcal{V}_{\ell,s}(t)|,
\label{eq:cell_occupancy}\\
B_{\ell,s}(t)&=\mathbb{I}[O_{\ell,s}(t)>0].
\label{eq:binary_occupancy}
\end{align}
For occupied cells, mean speed, total vehicle weight, and mean vehicle-category code are
\begin{align}
V_{\ell,s}(t)&=\frac{1}{O_{\ell,s}(t)}
\sum_{i\in\mathcal{V}_{\ell,s}(t)}v_i(t),
\label{eq:mean_vehicle_speed}\\
P_{\ell,s}(t)&=\sum_{i\in\mathcal{V}_{\ell,s}(t)}w_i,
\label{eq:cell_load}\\
C_{\ell,s}(t)&=\frac{1}{O_{\ell,s}(t)}
\sum_{i\in\mathcal{V}_{\ell,s}(t)}c_i.
\label{eq:mean_vehicle_type}
\end{align}
For an empty cell, $V_{\ell,s}=P_{\ell,s}=C_{\ell,s}=0$. Each cell is therefore represented by
\begin{equation}
\mathbf{Y}_{\ell,s}(t)=
[O_{\ell,s}(t),V_{\ell,s}(t),
P_{\ell,s}(t),C_{\ell,s}(t)]^T.
\label{eq:lane_cell_state}
\end{equation}
The complete traffic field has dimensions
$\mathbf{Y}(t)\in\mathbb{R}^{N_\ell\times N_x\times4}$, retaining the occupancy, speed, weight, and category channels.

\subsubsection{Spatiotemporal Load-Field Assembly}
\label{subsec:spatiotemporal_field}
The time-varying traffic tensor is
$\mathbfcal{Y}=\{\mathbf{Y}(t_1),\ldots,\mathbf{Y}(t_T)\}$.
Observed and future sequences are defined as
\begin{align}
\mathbf{X}_{\mathrm{obs}}&=
[\mathbf{Y}_{t-T_{\mathrm{obs}}+1},\ldots,\mathbf{Y}_t],
\label{eq:observed_sequence}\\
\mathbf{Y}_{\mathrm{future}}&=
[\mathbf{Y}_{t+1},\ldots,\mathbf{Y}_{t+T_{\mathrm{pred}}}].
\label{eq:future_sequence}
\end{align}
The implementation uses $T_{\mathrm{obs}}=20$ and $T_{\mathrm{pred}}=30$ frames. To accommodate differing units, the channels are normalized as
\begin{equation}
\widetilde{Y}^{(c)}=\frac{Y^{(c)}}{\alpha_c},
\label{eq:channel_normalization}
\end{equation}
where $\alpha_c$ denotes the scale for channel $c$. The weight channel is restored to physical units before bridge coupling. Fig.~\ref{fig:spatiotemporal_traffic} illustrates the evolution of occupancy, speed, and traffic loading.

\subsection{Physics-Guided Traffic Evolution Estimation}
\label{sec:physics_guided_evolution}

Primarily, the future bridge-traffic state is propagated by utilizing physically interpretable vehicle-motion models and is subsequently refined by a learned residual correction. The vehicle configuration at the final observed instant provides the initial condition for the future evolution. Longitudinal motion is described through the Intelligent Driver Model (IDM), while lane-changing behaviour is represented by utilizing a simplified MOBIL-like rule.

Let the complete vehicle configuration at the final observed instant be denoted by

\begin{equation}
\mathcal{V}_{t}
=
\left\{
z_i(t)
\right\}_{i=1}^{N_t}.
\end{equation}

Its subsequent physics-based evolution is represented by

\begin{equation}
\mathcal{V}_{t+\tau}^{\mathrm{phy}}
=
\Phi_{\mathrm{phy}}
\left(
\mathcal{V}_{t+\tau-1}^{\mathrm{phy}}
\right),
\qquad
\tau=1,\ldots,T_{\mathrm{pred}},
\end{equation}

here $\Phi_{\mathrm{phy}}$ represents the microscopic traffic-update operator. The resulting vehicle states are mapped back to the lane-space representation defined in Section~\ref{sec:traffic_state_representation}.

Fig. \ref{fig:spatiotemporal_traffic} illustrates a representative spatiotemporal traffic sequence in terms of occupancy, mean speed, and traffic-load distribution.

\begin{figure*}[!tbp]
\centering

\begin{minipage}[b]{0.30\textwidth}
    \centering
    \includegraphics[width=\linewidth]{3.1.png}\\
    (a)
\end{minipage}
\hfill
\begin{minipage}[b]{0.30\textwidth}
    \centering
    \includegraphics[width=\linewidth]{3.2.png}\\
    (b)
\end{minipage}
\hfill
\begin{minipage}[b]{0.30\textwidth}
    \centering
    \includegraphics[width=\linewidth]{3.3.png}\\
    (c)
\end{minipage}

\caption{Spatiotemporal bridge-traffic state representation for a representative traffic sequence: (a) vehicle occupancy, (b) mean vehicle speed, and (c) traffic-load distribution over the bridge span during the prediction interval.}
\label{fig:spatiotemporal_traffic}
\end{figure*}


\subsubsection{Longitudinal Vehicle Motion Using IDM}

For a following vehicle $i$ and its leader $j$, the longitudinal spacing can be written as:

\begin{equation}
d_i(t)
=
x_j(t)-x_i(t)-L_j,
\end{equation}

and the relative velocity is

\begin{equation}
\Delta v_i(t)
=
v_i(t)-v_j(t).
\end{equation}

The desired dynamic spacing is given by

\begin{equation}
d_i^{*}(t)
=
d_0
+
v_i(t)T_h
+
\frac{
v_i(t)\Delta v_i(t)
}{
2\sqrt{ab}
},
\end{equation}

where $d_0$ is the minimum spacing, $T_h$ the desired time headway, $a$ the maximum acceleration, and $b$ the comfortable deceleration.

The IDM acceleration is expressed as:

\begin{equation}
a_i(t)
=
a
\left[
1
-
\left(
\frac{v_i(t)}{v_0}
\right)^{\delta}
-
\left(
\frac{d_i^{*}(t)}{d_i(t)}
\right)^2
\right].
\end{equation}

The vehicle speed and position are then updated using the traffic time step $\Delta t$. This produces the longitudinal evolution required for reconstructing the future lane-space traffic field.


\subsubsection{Lane-Changing Behaviour Using a MOBIL-Like Rule}

Lane-changing behaviour is represented using a simplified MOBIL-like rule based on incentive and safety conditions. A candidate lane change is accepted when the expected acceleration benefit exceeds the prescribed threshold,

\begin{equation}
\Delta a_i
+
p
\left(
\Delta a_{\mathrm{new}}
+
\Delta a_{\mathrm{old}}
\right)
>
\Delta a_{\mathrm{th}},
\end{equation}

where $p$ is the politeness factor and $\Delta a_{\mathrm{th}}$ is the minimum incentive threshold.

The lane change must also satisfy the safety condition

\begin{equation}
a_{\mathrm{new}}^{\mathrm{after}}
\geq
-b_{\mathrm{safe}},
\end{equation}

where $b_{\mathrm{safe}}$ is the maximum admissible deceleration of the new following vehicle.

When both conditions are satisfied, the vehicle lane index is updated while its weight and vehicle category remain unchanged. The longitudinal position and speed continue to evolve according to the IDM update.


\subsubsection{Physics-Based Load-Field Rollout}

After each IDM and lane-change update, the propagated vehicle configuration is converted to the lane-space representation defined in Section~\ref{sec:traffic_state_representation}:

\begin{equation}
M_{t+\tau}^{\mathrm{phy}}
=
\mathcal{R}\left(
\mathcal{V}_{t+\tau}^{\mathrm{phy}}
\right),
\qquad
\tau=1,\ldots,T_{\mathrm{pred}},
\end{equation}

where $\mathcal{R}$ denotes the vehicle-state-to-load-field mapping operator.

Therefore, the complete physics-based future sequence is expressed as:

\begin{equation}
Y^{\mathrm{phy}}
=
\left[
M_{t+1}^{\mathrm{phy}},
M_{t+2}^{\mathrm{phy}},
\ldots,
M_{t+T_{\mathrm{pred}}}^{\mathrm{phy}}
\right].
\end{equation}

Each frame retains occupancy, speed, vehicle-weight, and vehicle-type information over the complete lane-space grid.


\subsubsection{Residual Spatiotemporal Correction}

The physics-based rollout provides an interpretable future traffic trajectory, however it may accumulate local errors in vehicle position, lane assignment, speed, and load intensity. Hence, a learned residual correction is applied to the physics-based sequence rather than predicting the complete future traffic field independently.

Let $X_t$ denote the observed traffic history and $Y^{\mathrm{phy}}$ the physics-based rollout. The residual correction is

\begin{equation}
\Delta\widehat{Y}
=
\mathcal{F}_{\theta}
\left(
X_t,
Y^{\mathrm{phy}}
\right),
\end{equation}

and the corrected traffic field is

\begin{equation}
\widehat{Y}
=
Y^{\mathrm{phy}}
+
\Delta\widehat{Y}.
\end{equation}

The correction module preserves the lane-space organization of the traffic field and estimates local deviations from the physically propagated trajectory. The resulting output has the same temporal, lane, spatial, and channel dimensions as the physics-based sequence.


\subsubsection{Occupancy Calibration and Training Objective}
\label{sec:occupancy_training}

The lane-space traffic representation is inherently sparse, because a large proportion of spatial cells remain unoccupied at a provided time instant. The count-valued channel $O_{\ell,s}$ preserves part of the traffic state, whereas the binary state $B_{\ell,s}$ is utilized for occupancy classification and load-support enforcement. Therefore, small non-zero predictions in reference-empty cells may establish artificial vehicle-load contributions when the predicted weight field is communicated to the VBI model. In the implemented residual model, the occupancy quantity entering the weighted binary cross-entropy loss is generated directly from the normalized predicted occupancy channel rather than from a separate occupancy-logit head. Let $\widehat{M}_{\tau,l,s}^{(o)}$ denote the normalized predicted occupancy channel and let $s_{\mathrm{occ}}$ denote its normalization factor. The probability deployed in the occupancy loss is

\begin{equation}
p_{\tau,l,s}^{\mathrm{occ}}
=
\operatorname{clip}
\left(
s_{\mathrm{occ}}\widehat{M}_{\tau,l,s}^{(o)},
10^{-5},
1-10^{-5}
\right).
\end{equation}



For validation calibration and downstream load-support enforcement, the normalized predicted occupancy channel is thresholded directly:
\begin{equation}
\widehat{B}_{\tau,l,s}
=
\mathbb{I}\left[
\widehat{M}_{\tau,l,s}^{(o)}
\geq \eta_{\mathrm{occ}}
\right].
\label{eq:predicted_binary_occupancy}
\end{equation}
The threshold $\eta_{\mathrm{occ}}$ is selected through the validation partition; the canonical residual model utilizes $\eta_{\mathrm{occ}}=0.23$. The predicted binary occupancy mask \(\hat B_{\tau,l,s}\) is used to suppress load-related predictions in cells classified as unoccupied so that artificial vehicle loads do not propagate into the downstream VBI analysis.

The residual correction model is trained by utilizing a composite objective designed to balance traffic-state reconstruction, consistency with the physics-based rollout, occupancy classification, and suppression of false activations. Let $\Omega$ denote the set of prediction indices over time, lane, longitudinal position, and traffic channel. The reconstruction term is expressed as

\begin{equation}
\mathcal{L}_{\mathrm{data}}
=
\frac{1}{|\Omega|}
\sum_{(\tau,l,s,c)\in\Omega}
\omega_c
\left(
\widehat{M}_{\tau,l,s,c}
-
M_{\tau,l,s,c}
\right)^2,
\end{equation}

where $\omega_c$ controls the relative contribution of the corresponding traffic channel.

To discourage unnecessary deviations from the physics-based evolution, a rollout-consistency term is introduced as

\begin{equation}
\mathcal{L}_{\mathrm{phy}}
=
\frac{1}{|\Omega|}
\sum_{(\tau,l,s,c)\in\Omega}
\left(
\widehat{M}_{\tau,l,s,c}
-
M_{\tau,l,s,c}^{\mathrm{phy}}
\right)^2.
\end{equation}

In the canonical implementation, both matrix-MSE terms are complementarily weighted by channel weights $(4,1,2,1)$. The data terminology utilizes an occupied-cell multiplier $1+8B_{\tau,l,s}$, while the Physics-Based consistency terminology utilizes the corresponding multiplier $1+B_{\tau,l,s}^{\mathrm{phy}}$.

As occupied cells form a smaller fraction of the traffic tensor, occupancy classification is further constrained through a weighted binary cross-entropy term, where $N$ represents the number of time-lane-cell occupancy entries comprised in the corresponding loss sum.

\begin{equation}
\begin{aligned}
\mathcal{L}_{\mathrm{occ}}
&=
-
\frac{1}{N}
\sum_{\tau,l,s}
\Bigg[
\lambda_{+}
B_{\tau,l,s}
\log
p_{\tau,l,s}^{\mathrm{occ}}
\\
&\qquad+
\left(
1-B_{\tau,l,s}
\right)
\log
\left(
1-p_{\tau,l,s}^{\mathrm{occ}}
\right)
\Bigg].
\end{aligned}
\end{equation}

where $B_{\tau,l,s}$ denotes the binary reference occupancy state defined by
Eq.~\eqref{eq:binary_occupancy}, and $\lambda_{+}$ compensates for class imbalance.


The executed empty-cell factor suppresses diffuse predicted occupancy specifically at locations that are empty in both the reference data and the Physics-Based rollout. Defining
\begin{equation}
E_{\tau,l,s}
=
\mathbb{I}\left[
B_{\tau,l,s}=0
\ \land\
B_{\tau,l,s}^{\mathrm{phy}}=0
\right],
\end{equation}
the corresponding penalty is
\begin{equation}
\mathcal{L}_{\mathrm{empty}}
=
\frac{
\sum_{\tau,l,s}
E_{\tau,l,s}
\left(
\widehat{M}_{\tau,l,s}^{(o)}
\right)^2
}{
\sum_{\tau,l,s}E_{\tau,l,s}
}.
\end{equation}

Since the vehicle-weight channel gives the direct interface with the subsequent VBI analysis, the implemented load-consistency term compares the total normalized vehicle weight at each future frame after summation over all lanes and longitudinal cells:
\begin{equation}
\mathcal{L}_{\mathrm{load}}
=
\frac{1}{T_{\mathrm{pred}}}
\sum_{\tau=1}^{T_{\mathrm{pred}}}
\left[
\sum_{l=1}^{L}\sum_{s=1}^{S}
\widehat{M}_{\tau,l,s}^{(w)}
-
\sum_{l=1}^{L}\sum_{s=1}^{S}
M_{\tau,l,s}^{(w)}
\right]^2.
\end{equation}
In addition, the implementation utilizes an $L_1$ penalty on the learned residual, $\mathcal{L}_{\mathrm{res}}=\operatorname{mean}|\Delta\widehat{Y}|$, and a temporal smoothness penalty, $\mathcal{L}_{\mathrm{smooth}}=\operatorname{mean}[(\widehat{M}_{\tau}-\widehat{M}_{\tau-1})^2]$.


Therefore, the implemented training objective is
\begin{equation}
\begin{aligned}
\mathcal{J}
={}&
\lambda_{\mathrm{data}}\mathcal{L}_{\mathrm{data}}
+
\lambda_{\mathrm{phy}}\mathcal{L}_{\mathrm{phy}}
+
\lambda_{\mathrm{occ}}\mathcal{L}_{\mathrm{occ}}
+
\lambda_{\mathrm{empty}}\mathcal{L}_{\mathrm{empty}}
\\
&+
\lambda_{\mathrm{res}}\mathcal{L}_{\mathrm{res}}
+
\lambda_{\mathrm{load}}\mathcal{L}_{\mathrm{load}}
+
\lambda_{\mathrm{smooth}}\mathcal{L}_{\mathrm{smooth}}.
\end{aligned}
\end{equation}
For the canonical run, the corresponding weights are $\lambda_{\mathrm{data}}=1.0$, $\lambda_{\mathrm{phy}}=0.05$, $\lambda_{\mathrm{occ}}=0.25$, $\lambda_{\mathrm{empty}}=1.2$, $\lambda_{\mathrm{res}}=0.01$, $\lambda_{\mathrm{load}}=0.10$, and $\lambda_{\mathrm{smooth}}=0.001$; the optional Markov regularization term is disabled. These settings remain fixed during held-out test evaluation.


The complete physics-guided traffic-load estimation procedure is summarized in Algorithm~\ref{alg:traffic_load_estimation}. The pseudocode propagates the vehicle set available at the end of the observation interval and does not specify a separate future-arrival mechanism.

\begin{algorithm}[!htbp]
\caption{Physics-Guided Spatiotemporal Load-Field Estimation}
\label{alg:traffic_load_estimation}
\footnotesize
\KwData{Observed traffic history $\mathbf{X}_t$, vehicle states
$\mathcal{V}_t$, horizon $T_{\mathrm{pred}}$, trained
$\mathcal{F}_\theta$, threshold $\eta_{\mathrm{occ}}$}
\KwResult{Occupancy-consistent future traffic field
$\widehat{\mathbf{Y}}$}
$\mathcal{V}^{\mathrm{phy}}_t\leftarrow\mathcal{V}_t$\;
\For{$\tau\leftarrow1$ \KwTo $T_{\mathrm{pred}}$}{
    \ForEach{vehicle $i$ in $\mathcal{V}^{\mathrm{phy}}_{t+\tau-1}$}{
        Identify the preceding vehicle in its lane\;
        Update position and speed using the IDM equations\;
        Apply the MOBIL incentive and safety rules to update the lane\;
        Preserve vehicle weight and category\;
    }
    $\mathcal{V}^{\mathrm{phy}}_{t+\tau}
    \leftarrow\{\text{updated vehicle states}\}$\;
    $\mathbf{M}^{\mathrm{phy}}_{t+\tau}
    \leftarrow\mathcal{R}(\mathcal{V}^{\mathrm{phy}}_{t+\tau})$\;
}
$\mathbf{Y}^{\mathrm{phy}}
\leftarrow[\mathbf{M}^{\mathrm{phy}}_{t+1},\ldots,
\mathbf{M}^{\mathrm{phy}}_{t+T_{\mathrm{pred}}}]$\;
$\widehat{\mathbf{Y}}\leftarrow\mathbf{Y}^{\mathrm{phy}}
+\mathcal{F}_\theta(\mathbf{X}_t,\mathbf{Y}^{\mathrm{phy}})$\;
\ForEach{prediction cell $(\tau,\ell,s)$}{
    $\widehat{B}_{\tau,\ell,s}\leftarrow
    \mathbb{I}[\widehat{M}^{(o)}_{\tau,\ell,s}
    \geq\eta_{\mathrm{occ}}]$\;
    Suppress the predicted vehicle load if $\widehat{B}_{\tau,\ell,s}=0$\;
}
$\widehat{\mathbf{Y}}\leftarrow
\operatorname{InverseScale}(\widehat{\mathbf{Y}})$\;
\Return{$\widehat{\mathbf{Y}}$}\;
\end{algorithm}


\subsection{Traffic-Bridge Coupling and Load Transfer}
\label{sec:traffic_bridge_transfer}

The estimated traffic-load field is defined over discrete lane--space cells, whereas the VBI formulation operates in structural coordinates. Therefore, a traffic-to-bridge coupling procedure is introduced to transfer the predicted vehicle loading to the structural model while retaining its spatial distribution. The vehicle-weight channel is first returned to its physical scale and preserved separately for each lane and longitudinal cell rather than being reduced to a single total bridge load.

Let $P_{\ell,s}(t)$ denote the vehicle load associated with lane $\ell$ and longitudinal cell $s$ at time $t$. Since the traffic and structural domains may have different physical lengths, the longitudinal cell location is first expressed in normalized form,

\begin{equation}
x_s^{\mathrm{tr}}
=
\left(s-\frac{1}{2}\right)
\frac{L_{\mathrm{tr}}}{S},
\qquad
s=1,\ldots,S,
\label{eq:traffic_cell_center}
\end{equation}

and the normalized longitudinal coordinate is

\begin{equation}
\xi_s
=
\frac{x_s^{\mathrm{tr}}}{L_{\mathrm{tr}}},
\label{eq:normalized_cell_coordinate}
\end{equation}

and mapped to the structural coordinate as

\begin{equation}
x_s^{\mathrm{b}}
=
\xi_s L_{\mathrm{b}},
\end{equation}

where $L_{\mathrm{tr}}$ and $L_{\mathrm{b}}$ denote the traffic-domain and structural-span lengths, respectively. The transformation preserves relative longitudinal position rather than equal physical distance. Thus, the 500~m traffic coordinate is normalized before it is projected onto the 318~m structural span.


\subsubsection{Full-Field Spatial Load Projection}

For lane $\ell$, the instantaneous total traffic load is acquired from the contributions of all longitudinal cells,

\begin{equation}
P_{\ell}(t)
=
\sum_{s=1}^{S}
P_{\ell,s}(t).
\end{equation}

Rather than applying this load at a single equivalent position, the contribution of each occupied lane-cell is retained in the structural projection. Let $\phi_r(x)$ denote the $r$-th bridge mode shape. Then, the modal contribution of the lane-space traffic field can be written in the general form as:

\begin{equation}
Q_r^{\mathrm{tr}}(t)
=
\sum_{\ell=1}^{L}
\sum_{s=1}^{S}
P_{\ell,s}(t)
\phi_r\left(x_s^{\mathrm{b}}\right).
\end{equation}

This quantity represents the modal projection of the spatial traffic-load field and is distinguished from the final VBI force generated by the coupled vehicle and bridge dynamics. Retaining the complete lane-cell distribution allows traffic states with similar total vehicle weight but distinct longitudinal load locations to produce distinct structural excitation.




\subsubsection{Observed-History Dynamic Initialization}

The structural state at the beginning of the future traffic interval depends on the loading that has already acted on the bridge during the observation period. Hence, the VBI system is not reinitialized from rest for each estimated traffic sequence. Instead, the common observed traffic history is first transferred through the same traffic-to-bridge coupling procedure and propagated through the VBI model.

For the observed traffic history $\mathbf{X}_{\mathrm{obs}}$, the dynamic state at the end of the observation interval is

\begin{equation}
\mathbf{z}_{0}
=
\mathcal{I}
\left(
\mathbf{X}_{\mathrm{obs}};
\boldsymbol{\theta}_{b},
\boldsymbol{\theta}_{v}
\right),
\end{equation}

where $\mathcal{I}(\cdot)$ denotes the dynamic initialization procedure and $\boldsymbol{\theta}_{b}$ and $\boldsymbol{\theta}_{v}$ contain the adopted bridge and vehicle parameters.

The initialized state contains the bridge modal and equivalent vehicle dynamic coordinates required to continue the coupled solution. The same state $\mathbf{z}_{0}$ is then used for all estimated future traffic-load fields. Consequently, differences among the resulting bridge responses arise from differences in the estimated traffic excitation rather than from inconsistent structural initial conditions.



\subsubsection{Multirate Traffic-to-Structural Time Transfer}

The traffic and structural models operate at different temporal resolutions: $\Delta t_{\mathrm{tr}}=0.5~\mathrm{s}$ for traffic states and $\Delta t_{\mathrm{s}}=0.02~\mathrm{s}$ for VBI integration.

For a structural time instant $t$ located between two consecutive traffic frames $t_k$ and $t_{k+1}$, the interpolation factor is

\begin{equation}
\alpha(t)
=
\frac{t-t_k}
{t_{k+1}-t_k}.
\end{equation}

The corresponding lane-cell load is evaluated as

\begin{equation}
P_{\ell,s}(t)
=
\left[1-\alpha(t)\right]
P_{\ell,s}(t_k)
+
\alpha(t)
P_{\ell,s}(t_{k+1}).
\end{equation}

Then, the interpolated loads are utilized to update the spatial participation terms required by the VBI solver. With the adopted time increments, each 0.5 s traffic interval is resolved using 25 structural integration increments.

Linear interpolation increases the temporal resolution of the excitation supplied to the structural integrator but does not reconstruct traffic-frequency content that is absent from the original 0.5 s traffic representation. The resulting continuously updated load field is supplied to the coupled vehicle-bridge equations described in the following section.








\subsection{VBI Dynamic Formulation}
\label{sec:vbi_dynamic_formulation}

The spatially and temporally resolved traffic loading from the traffic-to-bridge interface is introduced into the VBI formulation. The
bridge and vehicle subsystems interact through the relative vertical motion
between the vehicle suspension and the bridge surface. The continuous
formulation presented below provides the mechanical basis for the reduced-order
implementation developed in Section~\ref{sec:reduced_vbi}.


\subsubsection{Bridge and Vehicle Dynamic Representation}

The bridge is represented using an Euler-Bernoulli beam. Let $w(x,t)$ denote the vertical bridge displacement at longitudinal coordinate $x$. The governing equation is

\begin{equation}
\mu_b
\frac{\partial^2 w(x,t)}{\partial t^2}
+
EI
\frac{\partial^4 w(x,t)}{\partial x^4}
=
p(x,t),
\end{equation}

where $\mu_b$ is the bridge mass per unit length and $EI$ is the flexural rigidity.

For the simply supported representation,

\begin{equation}
w(0,t)=w(L_b,t)=0,
\end{equation}

and

\begin{equation}
\frac{\partial^2 w}{\partial x^2}(0,t)
=
\frac{\partial^2 w}{\partial x^2}(L_b,t)
=
0.
\end{equation}

The vehicle subsystem is represented by a vertical spring--damper
coordinate $z_i(t)$. For vehicle $i$, the relative suspension deformation is defined as

\begin{equation}
\delta_i(t)
=
z_i(t)
-
w\!\left(x_i(t),t\right)
-
r_i(t),
\end{equation}

where $x_i(t)$ is obtained from the traffic-motion model and $r_i(t)$ denotes the optional road-profile excitation.



The suspension relative velocity accounts for the
motion of the vehicle contact point along the bridge,
and is expressed as
\begin{equation}
\dot{\delta}_i(t)
=
\dot{z}_i(t)
-
\left[
\frac{\partial w}{\partial t}
+
\dot{x}_i(t)\frac{\partial w}{\partial x}
\right]_{x=x_i(t)}
-
\dot{r}_i(t).
\label{eq:relative_velocity}
\end{equation}




The corresponding vehicle-bridge suspension interaction force is expressed as


\begin{equation}
F_i(t)
=
k_i \delta_i(t)
+
c_i \dot{\delta}_i(t).
\end{equation}

Then, the vehicle vertical equation can be written as


\begin{equation}
m_i\ddot{z}_i(t)
+
c_i\dot{\delta}_i(t)
+
k_i\delta_i(t)
=
m_i g.
\label{eq:vehicle_motion_gravity}
\end{equation}

Here, the vertical displacement is defined as positive
downwards, and $m_i g$ denotes the gravitational
force acting on vehicle $i$, consistent with the adopted
coordinate convention.



For the reported numerical study, the road-profile term is set to zero,
as described in Section~\ref{sec:vbi_configuration}.

%\begin{figure*}[!tbp]
   % \centering
    %\includegraphics[width=\textwidth]{6.png}
 %   \caption{Residual architecture for physics-guided traffic-load prediction \& future bridge-load estimation.}
  %  \label{fig:arch}
%\end{figure*}

\subsubsection{Energy Basis of the Coupled System}

The bridge kinetic and bending strain energies are

\begin{equation}
T_b
=
\frac{1}{2}
\int_{0}^{L_b}\mu_b
\left(
\frac{\partial w(x,t)}{\partial t}
\right)^2
\,dx,
\end{equation}

and

\begin{equation}
U_b
=
\frac{1}{2}
\int_{0}^{L_b}
EI
\left(
\frac{\partial^2 w(x,t)}{\partial x^2}
\right)^2
\,dx.
\end{equation}

For the vertical vehicle coordinate,

\begin{equation}
T_{v,i}
=
\frac{1}{2}
m_i
\dot{z}_i^{\,2}(t),
\end{equation}

while the suspension energy is

\begin{equation}
U_{s,i}
=
\frac{1}{2}
k_i
\delta_i^{\,2}(t).
\end{equation}



The gravitational potential energy of the vehicle
subsystem is expressed as
\begin{equation}
U_g=-\sum_i m_i g z_i(t),
\end{equation}
where the downward-positive coordinate convention
is adopted.

The Lagrangian of the coupled system is therefore
\begin{equation}
\mathcal{L}
=
\left(T_b+\sum_i T_{v,i}\right)
-
\left(U_b+\sum_i U_{s,i}+U_g\right).
\label{eq:coupled_lagrangian}
\end{equation}

Applying Hamilton's principle with the virtual
work of non-conservative suspension damping yields
\begin{equation}
\delta\int_{t_1}^{t_2}\mathcal{L}\,dt
+
\int_{t_1}^{t_2}\delta W_{\mathrm{nc}}\,dt
=0,
\label{eq:hamilton_vbi}
\end{equation}
where $\delta W_{\mathrm{nc}}$ denotes the
virtual work contribution of the viscous
suspension forces.

The resulting continuous vehicle-bridge equations
provide the mechanical basis for the reduced-order
coupling approximation described in Section~II-F.




\subsubsection{Coupled Bridge Excitation}

For multiple vehicles acting on the bridge, the continuous bridge
equation may be written as

\begin{equation}
\begin{aligned}
\mu_b
\frac{\partial^2 w(x,t)}{\partial t^2}
+
EI
\frac{\partial^4 w(x,t)}{\partial x^4}
&=
\\[-0.2em]
&\quad
\sum_{i\in\mathcal{A}(t)}
F_i(t)
\delta_{\mathrm{D}}
\left(
x-x_i(t)
\right).
\end{aligned}
\end{equation}

where $\mathcal{A}(t)$ denotes the set of active vehicles and $\delta_{\mathrm{D}}(\cdot)$ is the Dirac delta function.

The continuous formulation is subsequently transformed into the reduced
modal representation described in Section~\ref{sec:reduced_vbi}, where the lane-space traffic
loading is retained during construction of the coupled excitation.

In the coupled formulation, the bridge receives the suspension
interaction forces transmitted through the vehicle coordinates.
The traffic-load field determines the effective lane loading and
the associated modal participation in the reduced-order model.
Accordingly, the direct traffic-load projection is not added
independently to the bridge force vector in addition to the
suspension-coupling forces, thereby avoiding double counting
of the same gravitational loading contribution.


\begin{algorithm}[!htbp]
\caption{Multirate Traffic-to-Structure VBI Response Computation}
\label{alg:vbi_response}
\scriptsize
\KwData{Observed and predicted load fields
$\mathbf{P}^{obs}$, $\widehat{\mathbf{P}}$; modal basis
$\phi_r$; parameters $\Theta_b,\Theta_v$;
steps $\Delta t_{tr},\Delta t_s$}
\KwResult{$w(x_p,t)$ for $x_p\in
\{L_b/4,L_b/2,3L_b/4\}$ and $\ddot w(L_b/2,t)$}
$N_{sub}\leftarrow\Delta t_{tr}/\Delta t_s$;
$\;\beta\leftarrow0.25,\ \gamma\leftarrow0.50$\;
$x_s^b\leftarrow(s-\frac{1}{2})L_b/S$,
$\ s=1,\ldots,S$\;
$\mathbf{P}^{all}\leftarrow[
\mathbf{P}^{obs}_1,\ldots,\mathbf{P}^{obs}_{T_{obs}},
\widehat{\mathbf{P}}_1,\ldots,\widehat{\mathbf{P}}_{T_{pred}}]$\;
Initialize $\mathbf{s}_{VBI}\leftarrow\mathbf{s}_{init}$\;
\For{$k\leftarrow1$ \KwTo $T_{obs}+T_{pred}-1$}{
    \For{$j\leftarrow1$ \KwTo $N_{sub}$}{
        $\rho\leftarrow j/N_{sub}$\;
        $\mathbf{P}^{*}\leftarrow
        (1-\rho)\mathbf{P}^{all}_{k}
        +\rho\mathbf{P}^{all}_{k+1}$\;
        \ForEach{lane $\ell$}{
            $P_\ell^{*}\leftarrow\sum_{s=1}^{S}P_{\ell,s}^{*}$;
            $\ m_\ell^{*}\leftarrow P_\ell^{*}/g$\;
            \eIf{$P_\ell^{*}>0$}{
                $h_{\ell r}^{*}\leftarrow
                \sum_{s=1}^{S}\frac{P_{\ell,s}^{*}}{P_\ell^{*}}
                \phi_r(x_s^b)$, for $r=1,\ldots,N_m$\;
            }{
                Retain previous $h_{\ell r}$ (or zero initially)\;
            }
        }
        $(\mathbf{M},\mathbf{C},\mathbf{K},\mathbf{f})
        \leftarrow\mathcal{A}(\mathbf{m}^{*},
        \mathbf{h}^{*};\Theta_b,\Theta_v)$\;
        $\mathbf{s}_{VBI}\leftarrow
        \operatorname{Newmark}(\mathbf{M},\mathbf{C},
        \mathbf{K},\mathbf{f},\mathbf{s}_{VBI},
        \Delta t_s,\beta,\gamma)$\;
        \If{$k\geq T_{obs}$}{
            Store future $\{q_r(t),\ddot q_r(t)\}_{r=1}^{N_m}$\;
        }
    }
}
\ForEach{$x_p\in\{L_b/4,L_b/2,3L_b/4\}$}{
    $w(x_p,t)\leftarrow\sum_{r=1}^{N_m}\phi_r(x_p)q_r(t)$\;
}
$\ddot w(L_b/2,t)\leftarrow
\sum_{r=1}^{N_m}\phi_r(L_b/2)\ddot q_r(t)$\;
\Return{$w(L_b/4,t),w(L_b/2,t),w(3L_b/4,t),
\ddot w(L_b/2,t)$}\;
\end{algorithm}

\subsection{Reduced-Order VBI Computational Formulation}
\label{sec:reduced_vbi}

The continuous vehicle-bridge
equations introduced in Section~\ref{sec:vbi_dynamic_formulation} are transformed into a reduced-order system
for numerical response computation. The bridge displacement is represented
using a finite modal basis, while the vehicle subsystem is described through
lane-equivalent vertical coordinates. The original lane-wise and longitudinal
traffic-load distribution is retained during construction of the modal
excitation so that the structural response remains sensitive to the spatial
location of the moving loads.

This formulation avoids
solving the continuous bridge equation directly at each structural time step
and provides an efficient representation for repeated VBI evaluation over the
complete traffic sequence.



\subsubsection{Modal Representation of the Bridge Response}

The bridge displacement
is approximated using $N_m$ retained vibration
modes,

\begin{equation}
w(x,t)
\approx
\sum_{r=1}^{N_m}
\phi_r(x)q_r(t),
\end{equation}

where $\phi_r(x)$ is the $r$-th modal basis
function and $q_r(t)$ is the corresponding
generalized coordinate.

For the adopted simply
supported basis,

\begin{equation}
\phi_r(x)
=
\sin
\left(
\frac{r\pi x}{L_b}
\right).
\end{equation}

The bridge velocity and
acceleration are obtained as

\begin{equation}
\dot{w}(x,t)
=
\sum_{r=1}^{N_m}
\phi_r(x)\dot{q}_r(t),
\end{equation}

and

\begin{equation}
\ddot{w}(x,t)
=
\sum_{r=1}^{N_m}
\phi_r(x)\ddot{q}_r(t).
\end{equation}

Projection onto the
modal basis gives

\begin{equation}
M_r\ddot{q}_r(t)
+
C_r\dot{q}_r(t)
+
K_rq_r(t)
=
Q_r(t),
\qquad
r=1,\ldots,N_m,
\end{equation}

where $M_r$, $C_r$, and $K_r$ are the generalized
modal mass, damping, and stiffness, respectively.

The total generalized modal force $Q_r(t)$ is distinguished from the
direct traffic-load projection. The latter is

\begin{equation}
Q_r^{\mathrm{tr}}(t)
=
\sum_{\ell=1}^{L}
\sum_{s=1}^{S}
P_{\ell,s}(t)
\phi_r
\left(
x_s^{\mathrm{b}}
\right).
\end{equation}

The direct projection therefore depends on the instantaneous longitudinal
distribution of traffic loading rather than only on the total bridge load.
The complete coupled-system forcing is represented through $\mathbf{f}(t)$
after vehicle--bridge interaction terms are assembled.

\begin{figure*}[!tbp]
\centering

% ==================== First row ====================

\begin{minipage}[t]{0.29\textwidth}
    \centering
    \includegraphics[width=\linewidth]{8.1.png}
\end{minipage}
\hfill
\begin{minipage}[t]{0.29\textwidth}
    \centering
    \includegraphics[width=\linewidth]{8.2.png}
\end{minipage}
\hfill
\begin{minipage}[t]{0.29\textwidth}
    \centering
    \includegraphics[width=\linewidth]{8.3.png}
\end{minipage}

\vspace{1.5mm}

% ==================== Second row ====================

\makebox[\textwidth][c]{%
    \begin{minipage}[t]{0.29\textwidth}
        \centering
        \includegraphics[width=\linewidth]{8.4.png}
    \end{minipage}
    \hspace{0.04\textwidth}
    \begin{minipage}[t]{0.29\textwidth}
        \centering
        \includegraphics[width=\linewidth]{8.5.png}
    \end{minipage}
}

\caption{Vehicle heterogeneity in the numerical dataset:
\textbf{(a)} Gross vehicle weight distribution;
\textbf{(b)} mean observed vehicle speed distribution;
\textbf{(c)} vehicle length distribution;
\textbf{(d)} maximum-acceleration parameter distribution;
\textbf{(e)} reaction-time parameter distribution.}

\label{fig:vehicle_heterogeneity}

\end{figure*}

\subsubsection{Coupled Vehicle-Bridge Reduced System}

The retained bridge
modal coordinates are coupled with lane-equivalent vertical vehicle
coordinates. Let

\begin{equation}
\mathbf{q}(t)
=
\begin{bmatrix}
q_1(t) &
q_2(t) &
\cdots &
q_{N_m}(t)
\end{bmatrix}^{T}
\end{equation}

denote the bridge modal
vector and

\begin{equation}
\mathbf{z}(t)
=
\begin{bmatrix}
z_1(t) &
z_2(t) &
\cdots &
z_L(t)
\end{bmatrix}^{T}
\end{equation}

the lane-equivalent vehicle coordinates. Each lane-equivalent coordinate is a reduced vertical degree of freedom used for coupling with the retained bridge modes; this representation does not establish an exact instantaneous aggregation of all individual vehicles in that lane.

The complete reduced
coordinate vector is

\begin{equation}
\mathbf{u}(t)
=
\begin{bmatrix}
\mathbf{q}(t) \\
\mathbf{z}(t)
\end{bmatrix}.
\end{equation}

For lane $\ell$, the bridge
deformation participating in the coupling is generated from the lane-specific
modal participation terms,

\begin{equation}
w_{\ell}(t)
=
\mathbf{h}_{\ell}^{T}(t)
\mathbf{q}(t),
\end{equation}

where $\mathbf{h}_{\ell}(t)$ is updated from the
interpolated lane-cell load distribution.

The relative vertical
deformation is then

\begin{equation}
\delta_{\ell}(t)
=
z_{\ell}(t)
-
\mathbf{h}_{\ell}^{T}(t)
\mathbf{q}(t),
\end{equation}

and the lane-equivalent
suspension force is written as

\begin{equation}
F_{\ell}(t)
=
k_{\ell}\delta_{\ell}(t)
+
c_{\ell}\dot{\delta}_{\ell}(t).
\end{equation}

The complete reduced
VBI system is assembled as

\begin{equation}
\mathbf{M}(t)\ddot{\mathbf{u}}(t)
+
\mathbf{C}(t)\dot{\mathbf{u}}(t)
+
\mathbf{K}(t)\mathbf{u}(t)
=
\mathbf{f}(t),
\end{equation}

where $\mathbf{M}(t)$, $\mathbf{C}(t)$, and $\mathbf{K}(t)$ are the assembled
system matrices and $\mathbf{f}(t)$ contains the
time-dependent traffic excitation.

For each lane $\ell$, the total traffic load and corresponding
instantaneous lane-equivalent mass are
\begin{equation}
P_{\ell}(t)=\sum_{s=1}^{S}P_{\ell,s}(t),\qquad
m_{\ell}(t)=\frac{P_{\ell}(t)}{g}.
\label{eq:lane_equiv_mass}
\end{equation}
For $P_{\ell}(t)>0$, the modal participation coefficient is
computed from the load-weighted cell positions:
\begin{equation}
h_{\ell r}(t)=\sum_{s=1}^{S}
\frac{P_{\ell,s}(t)}{P_{\ell}(t)}
\phi_r\!\left(x_s^{\mathrm{b}}\right).
\label{eq:lane_modal_participation}
\end{equation}
The equivalent lane suspension parameters are scaled as
\begin{equation}
\begin{aligned}
k_{\ell}(t)&=k_{\mathrm{ref}}
\frac{m_{\ell}(t)}{m_{\mathrm{ref}}},\\
c_{\ell}(t)&=c_{\mathrm{ref}}
\frac{m_{\ell}(t)}{m_{\mathrm{ref}}}.
\end{aligned}
\label{eq:lane_equiv_properties}
\end{equation}
Let $\mathbf{H}(t)\in\mathbb{R}^{L\times N_m}$ contain the
row vectors $\mathbf{h}_{\ell}^{T}(t)$, and let
$\mathbf{D}_k(t)=\operatorname{diag}(k_1,\ldots,k_L)$ and
$\mathbf{D}_c(t)=\operatorname{diag}(c_1,\ldots,c_L)$.
The assembled modal--vehicle mass matrix is
\begin{equation}
\mathbf{M}(t)=
\begin{bmatrix}
\mathbf{M}_b & \mathbf{0}\\
\mathbf{0} & \mathbf{M}_v(t)
\end{bmatrix},
\label{eq:assembled_mass}
\end{equation}
where $\mathbf{M}_b$ contains the bridge modal masses and
$\mathbf{M}_v(t)$ contains the lane-equivalent masses.
The assembled damping and stiffness matrices are

\begin{equation}
\scalebox{0.78}{$\displaystyle
\mathbf{C}(t)=
\begin{bmatrix}
\mathbf{C}_b+\mathbf{H}^{T}\mathbf{D}_c\mathbf{H}
&
-\mathbf{H}^{T}\mathbf{D}_c\\
-\mathbf{D}_c\mathbf{H}
&
\mathbf{D}_c
\end{bmatrix}
$}
\label{eq:assembled_damping}
\end{equation}

and

\begin{equation}
\scalebox{0.78}{$\displaystyle
\mathbf{K}(t)=
\begin{bmatrix}
\mathbf{K}_b+\mathbf{H}^{T}\mathbf{D}_k\mathbf{H}
&
-\mathbf{H}^{T}\mathbf{D}_k\\
-\mathbf{D}_k\mathbf{H}
&
\mathbf{D}_k
\end{bmatrix}
$}
\label{eq:assembled_stiffness}
\end{equation}

respectively. The adopted downward-positive force vector is
\begin{equation}
\mathbf{f}(t)=
\begin{bmatrix}
\mathbf{0}\\ g\mathbf{m}_v(t)
\end{bmatrix},
\label{eq:assembled_force}
\end{equation}
where $\mathbf{m}_v(t)$ is the column vector of active
lane-equivalent masses. Thus, gravity acts on the lane-equivalent
vehicle coordinates and reaches the bridge through the coupled
suspension terms; the direct projected traffic load is not added
again as a separate bridge force.

For a lane with negligible effective mass (below 50~kg),
the implemented numerical model uses a decoupled placeholder
mass of 50~kg and diagonal stiffness of 1~N/m, with zero
suspension coupling and gravitational forcing. Its modal
participation is therefore immaterial while the lane is inactive.


The coupling terms are updated at each structural time step using the evolving lane-space load distribution, allowing changes in the spatial traffic configuration to influence the bridge response continuously.


The implemented solver adopts an instantaneous lane-equivalent
coupling approximation. At each structural increment, the
interpolated traffic loads update the lane masses, load-weighted
modal participation, suspension parameters, and coupled matrices.
The resulting coefficients are treated as fixed within that
Newmark increment. In particular, the numerical suspension
velocity is approximated as
\begin{equation}
\dot{\delta}_{\ell}(t)\approx
\dot{z}_{\ell}(t)-
\mathbf{h}_{\ell}^{T}(t)\dot{\mathbf{q}}(t),
\label{eq:reduced_velocity_approx}
\end{equation}
without an explicit
$\dot{\mathbf{h}}_{\ell}^{T}(t)\mathbf{q}(t)$ contribution.
The model also does not separately considers momentum-exchange
terms associated with changing equivalent lane masses.
Consequently, the formulation is a reduced-order numerical
approximation, not an exact moving-contact multi-vehicle model.
Its physical accuracy is subject to these assumptions.


\subsubsection{Newmark-$\beta$ Time Integration of the Coupled System}

The reduced coupled
system is integrated using the Newmark average-acceleration scheme. At
structural time $t_n$,

\begin{equation}
\mathbf{M}_n\ddot{\mathbf{u}}_n
+
\mathbf{C}_n\dot{\mathbf{u}}_n
+
\mathbf{K}_n\mathbf{u}_n
=
\mathbf{f}_n.
\end{equation}

The numerical
parameters are

\begin{equation}
\beta = 0.25,
\qquad
\gamma = 0.50,
\end{equation}

with structural time
step

\begin{equation}
\Delta t_s = 0.02~\mathrm{s}.
\end{equation}

At each structural
increment, the interpolated lane-cell load field is utilized to update the
modal participation terms and the corresponding system matrices and excitation
vector. The resulting state is then advanced to the next time step using the
standard Newmark update relations.


\subsubsection{Bridge Response Reconstruction}

After solving the reduced coupled system, the physical bridge
response is reconstructed from the modal coordinates. The vertical displacement
is

\begin{equation}
u_b(\xi,t)
=
\sum_{r=1}^{N_m}
\phi_r(\xi)q_r(t),
\end{equation}

and the corresponding acceleration is

\begin{equation}
\ddot{u}_b(\xi,t)
=
\sum_{r=1}^{N_m}
\phi_r(\xi)\ddot{q}_r(t).
\end{equation}

In the present study, bridge displacement is evaluated at

\begin{equation}
\xi
=
\frac{L_b}{4},
\qquad
\frac{L_b}{2},
\qquad
\frac{3L_b}{4},
\end{equation}

while acceleration is evaluated at midspan,

\begin{equation}
\ddot{u}_{\mathrm{mid}}(t)
=
\ddot{u}_b
\left(
\frac{L_b}{2},t
\right).
\end{equation}

Subsequently, the reconstructed response histories are used for comparison with the simulation-reference bridge response.




\section{Experimental Setup}

\subsection{Bridge-Traffic Scenario and Load-Field Configuration}

The controlled numerical dataset comprises 250 independent sparse, normal, and dense IDM/MOBIL traffic scenarios, each with 121 frames sampled every 0.5~s. A 500~m, four-lane roadway is discretized into 50 longitudinal cells (10~m per cell), with four channels for occupancy, mean speed, total vehicle weight, and mean vehicle type. Passenger vehicles, buses/light trucks, rigid trucks, and heavy trailer-type vehicles are included. Since the Physics-Based forecast employs the traffic generator's IDM/MOBIL model family, the benchmark tests model-consistent synthetic conditions rather than field traffic or traffic-model mismatch. Each evaluation window contains 20 observed frames (10~s) and 30 future frames (15~s). Fig.~\ref{fig:vehicle_heterogeneity} illustrates vehicle heterogeneity, and Fig.~\ref{fig:traffic_regime} depicts traffic regimes and activities.

\subsection{Data Partitioning and Traffic-Estimation Configuration}

The 250 scenarios are split independently into 175 training, 38 validation, and 37 held-out test scenarios, preventing windows from the same realization from crossing partitions. The resulting sets contain 3000, 600, and 600 windows, respectively. Channel normalization is fixed across the splits, and predicted vehicle weights are returned to kN before structural analysis. The residual model uses AdamW, hidden dimension 96, batch size 32, and up to 40 epochs, with validation-based model selection and random seed 13. The training loss is defined in Section~\ref{sec:occupancy_training}, and all test hyperparameters remain fixed. Table~\ref{tab:dataset_configuration} summarizes the configuration.

\begin{table}[!htbp]
\centering
\caption{Dataset, split and forecast configuration}
\label{tab:dataset_configuration}
\footnotesize
\begin{tabular}{ll}
\toprule
\textbf{Item} & \textbf{Final configuration} \\
\midrule
Traffic scenarios       & 250 \\
Frames per scenario     & 121 \\
Traffic interval        & 0.5 s \\
Traffic-grid length     & 500 m \\
Lanes                   & 4 \\
Longitudinal cells      & 50 \\
Cell length             & 10 m \\
Channels                & Occupancy, speed, vehicle weight, \\
                        & and vehicle type \\
Observation window      & 20 frames / 10 s \\
Forecast horizon        & 30 frames / 15 s \\
Training scenarios      & 175 \\
Validation scenarios    & 38 \\
Test scenarios          & 37 \\
Training windows        & 3000 \\
Validation windows      & 600 \\
Test windows            & 600 \\
Canonical split seed    & 13 \\
Batch size              & 32 \\
Hidden dimension        & 96 \\
Maximum epochs          & 40 \\
\bottomrule
\end{tabular}
\end{table}




\subsection{VBI Numerical Model and Parameter Configuration}
\label{sec:vbi_configuration}

The structural-response
evaluation uses the reduced-order VBI formulation introduced in Section~\ref{sec:reduced_vbi}.
The structural span is 318~m, and the reduced model retains the first five modes. The adopted modal frequencies are 0.758, 1.533, 2.345, 3.506, and 4.022~Hz. The sine functions define the spatial modal basis, while these frequencies are treated as adopted reduced-order modal parameters. An effective bridge mass per unit length of 30,000 kg/m is
adopted as a numerical parameter for constructing the generalized modal masses.

The bridge damping
contribution and road-surface irregularity are both set to zero in the reported
numerical study. These terms are retained in the general formulation but are
not activated here, allowing differences in bridge response to be attributed primarily
to the traffic excitation and vehicle-bridge coupling.

The vehicle subsystem
is represented using lane-equivalent vertical coordinates coupled to the bridge
modes. The adopted reduced-order equivalent vehicle parameters are 63,294~kg for mass, 1,903,172~N/m for suspension stiffness, and 7,882~N$\cdot$s/m for damping. These fixed parameters define the reduced suspension representation used in the numerical model; they are not interpreted here as instantaneous aggregates of the vehicles occupying a lane.

The traffic-load field
is available at 0.5 s intervals, whereas the VBI system is integrated using a
structural time step of 0.02 s, giving 25 structural increments per traffic
interval. Lane--cell loads are linearly interpolated between consecutive traffic
states before modal projection. Time integration is performed using the Newmark
average-acceleration scheme with $\beta=0.25$ and $\gamma=0.50$.

The zero-damping, 0.02~s configuration is retained as the canonical baseline. Robustness is additionally evaluated using bridge damping ratios of 1\%, 2\%, and 5\%, and structural integration steps of 0.01 and 0.005~s. These sensitivity calculations use a deterministic subset of three windows from each of the 37 held-out scenarios (111 windows in total).

Before introducing the
estimated future traffic-load field, the common 20-frame observed traffic
history is propagated through the same VBI system to obtain the structural
state at the end of the observation interval. Then, this state is used as the
common initial condition for all subsequent 30-frame response evaluations.


\begin{figure*}[!tbp]
\centering
\begin{minipage}[b]{0.43\textwidth}
    \centering
    \includegraphics[width=\linewidth]{10.1.png}
\end{minipage}
\hfill
\begin{minipage}[b]{0.43\textwidth}
    \centering
    \includegraphics[width=\linewidth]{10.2.png}
\end{minipage}

\caption{Traffic-regime and activity distribution of the synthetic bridge-traffic
dataset. Panel (a) shows the number of sparse, normal, and dense scenarios,
while panel (b) reports the mean entered vehicles and lane-change events per
scenario for each traffic regime.}
\label{fig:traffic_regime}
\end{figure*}

\subsection{Comparative Methods and Evaluation Metrics}

All comparative methods
are evaluated using the same observed traffic history, traffic-to-bridge
transformation, reduced-order VBI model, modal basis, equivalent vehicle
parameters, and Newmark-$\beta$ integration settings. Only the estimated
traffic-load field supplied during the subsequent prediction interval is varied
among the methods.



\begin{table*}[!tbp]
\centering
\caption{Traffic-state and traffic-load estimation performance on the held-out test set.}
\label{tab:traffic_performance}
\scriptsize
\setlength{\tabcolsep}{4pt}
\renewcommand{\arraystretch}{1.15}

\resizebox{\textwidth}{!}{%
\begin{tabular}{lcccccc}
\toprule
\textbf{Method}
&
\begin{tabular}[c]{@{}c@{}}\textbf{Global}\\ \textbf{RMSE}\end{tabular}
&
\begin{tabular}[c]{@{}c@{}}\textbf{Occupied}\\ \textbf{RMSE}\end{tabular}
&
\textbf{F1}
&
\begin{tabular}[c]{@{}c@{}}\textbf{Empty}\\ \textbf{FPR}\end{tabular}
&
\begin{tabular}[c]{@{}c@{}}\textbf{Load RMSE}\\ \textbf{(kN)}\end{tabular}
&
\begin{tabular}[c]{@{}c@{}}\textbf{Total-load RMSE}\\ \textbf{(kN)}\end{tabular}
\\
\midrule

Persistence
& 0.1951 & 0.4433 & 0.1451 & 0.0901 & 72.70 & 706.54 \\

Physics-Based
& \textbf{0.1310} & 0.3361 & 0.5182
& \textbf{0.0298} & \textbf{45.23} & 917.54 \\

ConvLSTM
& 0.2683 & 0.2752 & 0.1886 & 1.0000 & 58.39 & 852.99 \\

Transformer
& 0.2186 & 0.3246 & 0.2115 & 0.7442 & 54.06 & 671.57 \\

Direct Prediction
& 0.2208 & 0.3240 & 0.2061 & 0.7843 & 53.60 & 576.99 \\

Residual Correction
& 0.1568 & \textbf{0.2306} & \textbf{0.5370}
& 0.0334 & 45.30 & \textbf{512.02} \\

\bottomrule
\end{tabular}%
}
\end{table*}



\subsubsection{Comparative Traffic-Estimation Methods}

The six
traffic-estimation strategies are evaluated.

\textbf{Persistence:}
The final observed traffic state is repeated throughout the 15 s estimation
interval. This provides a simple baseline without explicit future
vehicle-motion modelling.

\textbf{Physics-Based:}
The final observed vehicle configuration is propagated using the IDM-based
longitudinal-motion model and the simplified MOBIL-like lane-change rule
described in Section~\ref{sec:physics_guided_evolution}. The resulting vehicle states are subsequently mapped
to the lane--space traffic-load field without any learned correction.

\textbf{ConvLSTM:}
A convolutional recurrent baseline is deployed to estimate the future traffic
tensor directly from the observed sequence while preserving the spatial
organization of the lane-space grid.

\textbf{Transformer:}
A data-driven attention-based sequence model is utilized to estimate the
subsequent lane-space traffic field from the observed history without explicit
car-following or lane-changing constraints.

\textbf{Direct Prediction:}
This model predicts the complete future traffic field directly from the
observed sequence and does not use the Physics-Based rollout as an additive
reference.

\textbf{Residual Correction:}
The residual model uses the Physics-Based rollout as its reference trajectory
and estimates only the remaining spatiotemporal deviation. Thus, the learned
correction is applied to the physically propagated traffic field rather than
replacing it completely.

For traffic-level
evaluation, the continuous Residual Correction output is retained. Before its
vehicle-weight channel is supplied to the VBI model, an occupancy-based support
constraint is applied using the validation-selected threshold
$\eta_{\mathrm{occ}}=0.23$:

\begin{equation}
\widehat{P}^{\,\mathrm{SC}}_{\tau,l,s}
=
\widehat{B}_{\tau,l,s}\widehat{P}_{\tau,l,s}.
\label{eq:support_consistent_load}
\end{equation}

Here, $\widehat{B}_{\tau,l,s}$ is the predicted binary occupancy-support mask
defined in Eq.~\eqref{eq:predicted_binary_occupancy}.
The resulting
structural-input case is referred to as Support-Consistent Residual Correction.
It is not an independently trained model. It is the Residual Correction output after enforcing consistency between predicted occupancy and vehicle-load support.

\textbf{Load-Matched Raw control:}
To separate the effect of spatial support correction from that of aggregate load reduction, a frame-wise load-matched control is constructed from the Raw Residual field. For each prediction step,
\begin{equation}
\alpha_{\tau}
=
\frac{
\sum_{l=1}^{L}\sum_{s=1}^{S}
\widehat{P}^{\,\mathrm{SC}}_{\tau,l,s}
}{
\sum_{l=1}^{L}\sum_{s=1}^{S}
\widehat{P}^{\,\mathrm{raw}}_{\tau,l,s}
},
\qquad
\widehat{P}^{\,\mathrm{LM}}_{\tau,l,s}
=
\alpha_{\tau}
\widehat{P}^{\,\mathrm{raw}}_{\tau,l,s}.
\label{eq:load_matched_control}
\end{equation}
Thus, the Load-Matched Raw field retains the spatial distribution of the Raw Residual prediction while matching the Support-Consistent total load at every future frame.

Because multiple temporal windows originate from the same traffic scenario, uncertainty in the support-consistency comparisons is evaluated at the scenario level. Window-level errors are first averaged within each of the 37 independent test scenarios, and 95\% confidence intervals for paired mean differences are obtained using 20,000 nonparametric bootstrap resamples of the scenario clusters.



\subsubsection{Traffic-State and Traffic-Load Evaluation Metrics}

Traffic-estimation
performance is evaluated using complementary matrix-level, occupancy-based, and
load-based metrics.

\textbf{Global normalized RMSE:}

For the metric definitions below, $N_c=4$ denotes the number of traffic channels. The global RMSE is computed over all prediction steps, lanes, longitudinal cells, and traffic channels:

\begin{equation}
\begin{aligned}
\mathrm{RMSE}_{\mathrm{global}}
&=
\Bigg[
\frac{1}
{T_{\mathrm{pred}}LSN_c}
\sum_{\tau=1}^{T_{\mathrm{pred}}}
\sum_{l=1}^{L}
\\
&\qquad\qquad
\sum_{s=1}^{S}
\sum_{c=1}^{N_c}
\left(
\widehat{M}_{\tau,l,s,c}
-
M_{\tau,l,s,c}
\right)^2
\Bigg]^{1/2}.
\end{aligned}
\end{equation}

In this matrix-level metric, $M$ and $\widehat{M}$ denote the channel-scaled traffic tensors defined by the normalization procedure above. This metric represents the overall traffic-tensor reconstruction error.

\textbf{Occupied-cell RMSE:}

Since the traffic field is sparse, the prediction error is also evaluated only
over reference-occupied cells:

\begin{equation}
\mathrm{RMSE}_{\mathrm{occ}}
=
\sqrt{
\frac{1}{|\Omega_{\mathrm{occ}}|}
\sum_{(\tau,l,s,c)\in\Omega_{\mathrm{occ}}}
\left(
\widehat{M}_{\tau,l,s,c}
-
M_{\tau,l,s,c}
\right)^2
},
\end{equation}

where

\begin{equation}
\Omega_{\mathrm{occ}}
=
\left\{
(\tau,l,s,c)
:
B_{\tau,l,s}=1
\right\}.
\end{equation}

This metric emphasizes
prediction quality in regions containing actual vehicle information.

\textbf{Occupancy F1-score:}

Occupancy classification performance is evaluated using

\begin{equation}
F_1
=
\frac{2TP}
{2TP+FP+FN},
\end{equation}

where $TP$, $FP$, and $FN$ denote true-positive, false-positive, and false-negative occupancy classifications, respectively. These counts compare the binary reference state $B_{\tau,l,s}$ with the predicted state $\widehat{B}_{\tau,l,s}$.

\textbf{Empty-cell false-positive rate:}

False activations within reference-empty cells are measured using the
empty-cell false-positive rate. This metric is particularly relevant for VBI
analysis, since false occupancy may introduce artificial vehicle-load
contributions.

\textbf{Lane-cell traffic-load RMSE:}

The vehicle-weight channel is first transformed back to its physical unit of
kN. The lane-cell load RMSE is then evaluated directly from the reference and
estimated vehicle-load fields:

\begin{equation}
\begin{aligned}
\mathrm{RMSE}_{\mathrm{load}}
&=
\Bigg[
\frac{1}
{T_{\mathrm{pred}}LS}
\sum_{\tau=1}^{T_{\mathrm{pred}}}
\sum_{l=1}^{L}
\\
&\qquad\qquad
\sum_{s=1}^{S}
\left(
\widehat{P}_{\tau,l,s}
-
P_{\tau,l,s}
\right)^2
\Bigg]^{1/2}.
\end{aligned}
\end{equation}

\textbf{Total-load RMSE:}

At
each prediction step, the total bridge traffic load is evolved by summing the
vehicle-weight field over all lanes and longitudinal cells. The corresponding
RMSE is

\begin{equation}
\begin{aligned}
\mathrm{RMSE}_{\mathrm{total}}
&=
\Bigg[
\frac{1}{T_{\mathrm{pred}}}
\sum_{\tau=1}^{T_{\mathrm{pred}}}
\Bigg(
\sum_{l=1}^{L}
\sum_{s=1}^{S}
\widehat{P}_{\tau,l,s}
\\
&\qquad\qquad
-
\sum_{l=1}^{L}
\sum_{s=1}^{S}
P_{\tau,l,s}
\Bigg)^2
\Bigg]^{1/2}.
\end{aligned}
\end{equation}

Both lane-cell and
total-load errors are reported in kN. Using both local and aggregate load
metrics is important, since similar total traffic loads can still correspond to
distinct spatial load distributions over the bridge.



\subsubsection{Bridge-Response Evaluation Metrics}

Bridge-response
accuracy is evaluated by comparing the response generated from each estimated
traffic-load field with the simulation-reference response generated using the
same VBI model.

The
vertical displacement response is evaluated at the quarter-span, midspan, and
three-quarter-span locations. For a selected bridge location $\xi$, the displacement RMSE
is

\begin{equation}
\mathrm{RMSE}_{u}(\xi)
=
\sqrt{
\frac{1}{N_t}
\sum_{n=1}^{N_t}
\left[
\widehat{u}(\xi,t_n)
-
u^{\mathrm{ref}}(\xi,t_n)
\right]^2
}.
\end{equation}

Agreement
of the displacement histories is additionally quantified using the coefficient
of determination,

\begin{equation}
R_u^2(\xi)
=
1-
\frac{
\displaystyle
\sum_{n=1}^{N_t}
\left[
u^{\mathrm{ref}}(\xi,t_n)
-
\widehat{u}(\xi,t_n)
\right]^2
}{
\displaystyle
\sum_{n=1}^{N_t}
\left[
u^{\mathrm{ref}}(\xi,t_n)
-
\overline{u}^{\,\mathrm{ref}}(\xi)
\right]^2
}.
\end{equation}

The
bridge acceleration response is evaluated at midspan. Its RMSE is

\begin{equation}
\begin{aligned}
\mathrm{RMSE}_{a}
&=
\Bigg[
\frac{1}{N_t}
\sum_{n=1}^{N_t}
\Bigg(
\widehat{\ddot{u}}
\left(
\frac{L_b}{2},t_n
\right)
\\
&\qquad\qquad
-
\ddot{u}^{\mathrm{ref}}
\left(
\frac{L_b}{2},t_n
\right)
\Bigg)^2
\Bigg]^{1/2}.
\end{aligned}
\end{equation}

and
the corresponding $R^2$ value is calculated in
the same manner as for displacement.

To
assess peak structural demand, the absolute peak midspan displacement is
computed for each test sequence, and the corresponding mean absolute error is
reported as

\begin{equation}
\begin{aligned}
\mathrm{MAE}_{\mathrm{peak}}
&=
\frac{1}{N_{\mathrm{test}}}
\sum_{j=1}^{N_{\mathrm{test}}}
\Bigg|
\max_{n}
\left|
\widehat{u}_{j}
\left(
\frac{L_b}{2},t_n
\right)
\right|
\\
&\qquad
-
\max_{n}
\left|
u_{j}^{\mathrm{ref}}
\left(
\frac{L_b}{2},t_n
\right)
\right|
\Bigg|.
\end{aligned}
\end{equation}

Displacement
errors are reported in mm, acceleration errors in m/s$^2$, and load-related errors
in kN.

For all structural-response comparisons, the same VBI solver, structural parameters, numerical integration settings, and observed-history initialization are employed. The simulation-reference response is obtained by applying the reference traffic-load field to the same reduced-order VBI model, ensuring that differences between methods arise from their estimated future loading.





\section{Results and Discussion}

\subsection{Traffic-State and Traffic-Load Estimation Performance}
\label{sec:traffic_results}

Table~\ref{tab:traffic_performance} reports performance over 600 held-out windows. Physics-Based forecasting yields the lowest global normalized RMSE (0.1310), empty-cell false-positive rate (0.0298), and lane-cell load RMSE (45.23~kN). Residual Correction achieves the lowest occupied-cell RMSE (0.2306), highest occupancy F1 (0.5370), and lowest aggregate total-load RMSE (512.02~kN), compared with 917.54~kN for Physics-Based forecasting. The purely data-driven baselines show that a lower total-load error can coexist with poorer occupancy support, especially for ConvLSTM, Transformer, and Direct Prediction. Thus, aggregate loading, local traffic reconstruction, and occupancy consistency must be evaluated separately.

Fig.~\ref{fig:horizon_load_rmse} shows horizon-wise total-load error, and Fig.~\ref{fig:spatiotemporal_load_fields} compares the estimated spatial load fields. Their structural-response consequences are examined next.

\begin{figure}[!htbp]
\centering
\includegraphics[width=0.94\columnwidth]{12.png}
\caption{Horizon-wise total traffic-load RMSE for the evaluated forecasting methods.}
\label{fig:horizon_load_rmse}
\end{figure}



\begin{figure*}[!tbp]
\centering
% Balanced, centered two-by-two layout for the four traffic-load fields.
\noindent\makebox[\textwidth][c]{%
\begin{minipage}[t]{0.41\textwidth}
\centering
\includegraphics[width=\linewidth]{14.1.png}
\end{minipage}%
\hspace{0.04\textwidth}%
\begin{minipage}[t]{0.41\textwidth}
\centering
\includegraphics[width=\linewidth]{14.2.png}
\end{minipage}%
}
\par\vspace{1mm}
\noindent\makebox[\textwidth][c]{%
\begin{minipage}[t]{0.41\textwidth}
\centering
\includegraphics[width=\linewidth]{14.3.png}
\end{minipage}%
\hspace{0.04\textwidth}%
\begin{minipage}[t]{0.41\textwidth}
\centering
\includegraphics[width=\linewidth]{14.4.png}
\end{minipage}%
}
\par
\caption{Spatiotemporal lane-summed traffic-load fields for a common held-out traffic
sequence: \textbf{(a)} simulation ground truth,
\textbf{(b)} physics-based traffic evolution,
\textbf{(c)} raw residual correction, and
\textbf{(d)} support-consistent residual correction.}
\label{fig:spatiotemporal_load_fields}
\end{figure*}








\subsection{Bridge Dynamic Response under Estimated Traffic Loads}
\label{sec:bridge_response_results}

All estimated traffic fields are evaluated with identical VBI parameters, structural integration, and observed-history initialization. Table~\ref{tab:bridge_response_performance} compares six primary structural inputs; the Raw Residual input is analyzed separately as an occupancy-support diagnostic. Physics-Based forecasting gives the lowest midspan displacement RMSE (3.427~mm) and acceleration RMSE (0.0269~m/s$^2$), while Persistence has the lowest quarter-span displacement RMSE (3.190~mm). Support-Consistent Residual yields 3.727~mm and 0.0321~m/s$^2$ at midspan, respectively. Purely data-driven forecasts have larger bridge-response errors, notably in acceleration, despite their occasionally favorable traffic-load metrics.

\begin{table*}[!tbp]
\centering
\caption{Bridge-response performance obtained from the evaluated traffic-load inputs.}
\label{tab:bridge_response_performance}
\scriptsize
\setlength{\tabcolsep}{4pt}
\renewcommand{\arraystretch}{1.15}
\resizebox{\textwidth}{!}{%
\begin{tabular}{lccccccc}
\toprule
\textbf{Traffic input}
& \begin{tabular}[c]{@{}c@{}}$L_b/4$ \textbf{RMSE}\\ \textbf{(mm)}\end{tabular}
& \begin{tabular}[c]{@{}c@{}}$L_b/2$ \textbf{RMSE}\\ \textbf{(mm)}\end{tabular}
& \begin{tabular}[c]{@{}c@{}}$3L_b/4$ \textbf{RMSE}\\ \textbf{(mm)}\end{tabular}
& \begin{tabular}[c]{@{}c@{}}\textbf{Acc. RMSE}\\ \textbf{(m/s$^2$)}\end{tabular}
& \begin{tabular}[c]{@{}c@{}}\textbf{Peak MAE}\\ \textbf{(mm)}\end{tabular}
& \begin{tabular}[c]{@{}c@{}}\textbf{Disp.}\\ $\mathbf{R^2}$\end{tabular}
& \begin{tabular}[c]{@{}c@{}}\textbf{Acc.}\\ $\mathbf{R^2}$\end{tabular}
\\
\midrule
Persistence
& \textbf{3.190} & 4.392 & 3.416 & 0.0281 & 4.416 & 0.571 & 0.364 \\
Physics-Based
& 3.894 & \textbf{3.427} & \textbf{1.856} & \textbf{0.0269} & \textbf{2.525} & \textbf{0.765} & \textbf{0.412} \\
ConvLSTM
& 5.391 & 7.236 & 5.238 & 0.1006 & 6.423 & -0.825 & -15.517 \\
Transformer
& 4.061 & 5.160 & 3.790 & 0.0792 & 3.517 & 0.523 & -1.838 \\
Direct Prediction
& 4.011 & 5.136 & 3.675 & 0.0869 & 4.002 & 0.517 & -2.545 \\
Support-Consistent Residual
& 4.062 & 3.727 & 2.101 & 0.0321 & 2.530 & 0.730 & 0.350 \\
\bottomrule
\end{tabular}%
}
\end{table*}





Response histories and scenario-level error distributions are shown in Figs.~\ref{fig:representative_bridge_response} and~\ref{fig:bridge_error_distribution}, respectively. The differing rankings of traffic-level and structural-response accuracy are discussed in Section~\ref{sec:load_response_fidelity}.

\begin{figure*}[!tbp]
\centering

\begin{minipage}[t]{0.43\textwidth}
    \centering
    \includegraphics[width=\linewidth]{15.1.png}
\end{minipage}
\hfill
\begin{minipage}[t]{0.43\textwidth}
    \centering
    \includegraphics[width=\linewidth]{15.2.png}
\end{minipage}

\caption{Representative displacement and acceleration history.}
\label{fig:representative_bridge_response}
\end{figure*}



\begin{figure*}[!tbp]
\centering

% ---------- First row ----------
\begin{minipage}[t]{0.29\textwidth}
    \centering
    \includegraphics[width=\linewidth]{18.1.png}
\end{minipage}
\hfill
\begin{minipage}[t]{0.29\textwidth}
    \centering
    \includegraphics[width=\linewidth]{18.2.png}
\end{minipage}
\hfill
\begin{minipage}[t]{0.29\textwidth}
    \centering
    \includegraphics[width=\linewidth]{18.3.png}
\end{minipage}

\vspace{1.5mm}

% ---------- Second row: centered ----------
\makebox[\textwidth][c]{%
    \begin{minipage}[t]{0.29\textwidth}
        \centering
        \includegraphics[width=\linewidth]{18.4.png}
    \end{minipage}
    \hspace{0.02\textwidth}
    \begin{minipage}[t]{0.29\textwidth}
        \centering
        \includegraphics[width=\linewidth]{18.5.png}
    \end{minipage}
}

\caption{Variation of bridge-response errors across scenarios
\textbf{(a)} $L_b/4$ displacement error distribution;
\textbf{(b)} $L_b/2$ displacement error distribution;
\textbf{(c)} $3L_b/4$ displacement error distribution;
\textbf{(d)} midspan acceleration error distribution;
\textbf{(e)} peak-midspan-displacement error distribution.}
\label{fig:bridge_error_distribution}
\end{figure*}





\subsection{Traffic-Load Estimation versus Bridge-Response Fidelity}
\label{sec:load_response_fidelity}

The results in Sections~\ref{sec:traffic_results} and~\ref{sec:bridge_response_results} show that lower traffic-load estimation error does not necessarily
lead to more accurate bridge-response prediction. Residual Correction achieves
the lowest total-load RMSE of 512.02 kN, whereas the Physics-Based method gives
a substantially larger value of 917.54 kN. However, when the corresponding
unmodified traffic fields are introduced into the VBI model, the Physics-Based
method produces lower midspan displacement and acceleration RMSE values of
3.427 mm and 0.0269 m/s$^2$, compared with 4.219 mm and 0.0552 m/s$^2$ for the raw
residual estimate.

This difference arises since
the structural response depends not only on the total vehicle load but also on
its spatial and temporal distribution over the bridge. Hence, two traffic
estimates with similar or improved aggregate loading can produce different
bridge responses when the vehicle loads are positioned distinctly along the
span.

Vehicle-level trajectory analysis further shows that, on average, 29.3\% of the vehicle weight occurring during the future prediction interval is associated with vehicles that are absent at the final observed traffic state. This future-arrival weight fraction is positively associated with the Physics-Based total-load RMSE (Pearson $r=0.471$; Spearman $\rho=0.559$). Since these vehicles are unavailable when the deterministic rollout is initialized, this association is consistent with unobserved future arrivals being an important source of Physics-Based forecast mismatch. The median RMSE between trajectory-derived and matrix-derived reference total load is only 0.168~kN, indicating close consistency between the two reference representations.

These results demonstrate that traffic-load accuracy and bridge-response fidelity should be evaluated separately. Fig.~\ref{fig:traffic_bridge_error_relationship} shows the corresponding relationships between traffic-load errors and structural-response errors. The influence of occupancy-load support consistency on this
relationship is examined in the following section.

\begin{figure*}[!tbp]
\centering

% ---------- First row ----------
\begin{minipage}[t]{0.29\textwidth}
    \centering
    \includegraphics[width=\linewidth]{19.1.png}
\end{minipage}
\hfill
\begin{minipage}[t]{0.29\textwidth}
    \centering
    \includegraphics[width=\linewidth]{19.2.png}
\end{minipage}
\hfill
\begin{minipage}[t]{0.29\textwidth}
    \centering
    \includegraphics[width=\linewidth]{19.3.png}
\end{minipage}

\vspace{1.5mm}

% ---------- Second row: centered ----------
\makebox[\textwidth][c]{%
    \begin{minipage}[t]{0.29\textwidth}
        \centering
        \includegraphics[width=\linewidth]{19.4.png}
    \end{minipage}
    \hspace{0.02\textwidth}
    \begin{minipage}[t]{0.29\textwidth}
        \centering
        \includegraphics[width=\linewidth]{19.5.png}
    \end{minipage}
}

\caption{Relationship between traffic-load and bridge-response errors
\textbf{(a)} Total-load error versus midspan displacement error;
\textbf{(b)} total-load error versus acceleration error;
\textbf{(c)} lane-cell load error versus displacement error;
\textbf{(d)} lane-cell load error versus acceleration error;
\textbf{(e)} total-load error versus peak-displacement error.}
\label{fig:traffic_bridge_error_relationship}
\end{figure*}




\subsection{Effect of Occupancy-Load Support Consistency}
\label{sec:support_consistency_results}

The raw residual estimate
may assign non-zero vehicle loads to cells predicted as unoccupied. To reduce this inconsistency, the vehicle-load field is masked using the occupancy
threshold selected from the validation data before it is supplied to the VBI model.
This produces the Support-Consistent Residual input without retraining the
forecasting model.

Applying this constraint substantially reduces the extreme structural response, as shown in Fig.~\ref{fig:support_consistency_effect}. The maximum midspan displacement decreases from 191.12 to 40.61 mm, while the maximum midspan acceleration decreases from 5.281 to 0.605 m/s$^2$. These maximum amplitudes describe extreme predicted responses and are distinct from the sequence-averaged RMSE and MAE values reported in Table~\ref{tab:bridge_response_performance}. These results indicate that
maintaining consistency between predicted occupancy and vehicle-load location
is important when estimated traffic fields are utilized for bridge dynamic
analysis.

To verify that this improvement is not explained only by the reduction in aggregate predicted load, the Load-Matched Raw control in Eq.~\eqref{eq:load_matched_control} was evaluated over all 600 test windows. The maximum frame-wise difference between its total load and the Support-Consistent total load was $9.1\times10^{-13}$~kN, and no frame required an undefined zero-load rescaling. Table~\ref{tab:load_matched_control} summarizes the resulting response errors.

\begin{table}[!htbp]
\centering
\caption{Load-matched support-consistency control over the 600 held-out test windows.}
\label{tab:load_matched_control}
\scriptsize
\setlength{\tabcolsep}{3pt}
\renewcommand{\arraystretch}{1.12}
\begin{tabular}{lccc}
\toprule
\textbf{Traffic input}
& \begin{tabular}[c]{@{}c@{}}\textbf{Midspan disp.}\\\textbf{RMSE (mm)}\end{tabular}
& \begin{tabular}[c]{@{}c@{}}\textbf{Midspan acc.}\\\textbf{RMSE (m/s$^2$)}\end{tabular}
& \begin{tabular}[c]{@{}c@{}}\textbf{Peak-midspan}\\\textbf{error (mm)}\end{tabular}
\\
\midrule
Raw Residual & 4.219 & 0.0552 & 5.002 \\
Load-Matched Raw & 4.395 & 0.0439 & 3.760 \\
Support-Consistent Residual & 3.727 & 0.0321 & 2.955 \\
Physics-Based & 3.427 & 0.0269 & 2.966 \\
\bottomrule
\end{tabular}
\end{table}

At the scenario level, Support-Consistent Residual improves over Load-Matched Raw by 0.654~mm in midspan displacement RMSE (95\% CI: 0.412--1.043~mm), 0.0114~m/s$^2$ in acceleration RMSE (95\% CI: 0.0053--0.0216~m/s$^2$), and 0.754~mm in peak-midspan displacement error (95\% CI: 0.280--1.535~mm). Improvements occur in 35/37, 33/37, and 30/37 scenarios, respectively. Hence, the structural benefit of occupancy-load support consistency cannot be attributed solely to a reduction in total predicted load; the spatial support of the predicted loading has an independent effect on the reconstructed bridge response. For the direct Raw Residual versus Support-Consistent comparison, the scenario-level confidence interval for midspan displacement RMSE includes zero, whereas those for acceleration and peak-midspan displacement error remain positive; the displacement benefit should therefore be interpreted as more scenario dependent.


\begin{figure*}[!tbp]
\centering

% ---------- First row ----------
\begin{minipage}[t]{0.30\textwidth}
    \centering
    \includegraphics[width=\linewidth]{20.1.png}
\end{minipage}
\hfill
\begin{minipage}[t]{0.30\textwidth}
    \centering
    \includegraphics[width=\linewidth]{20.2.png}
\end{minipage}
\hfill
\begin{minipage}[t]{0.30\textwidth}
    \centering
    \includegraphics[width=\linewidth]{20.3.png}
\end{minipage}

\vspace{1.5mm}

% ---------- Second row: centered ----------
\makebox[\textwidth][c]{%
    \begin{minipage}[t]{0.30\textwidth}
        \centering
        \includegraphics[width=\linewidth]{20.4.png}
    \end{minipage}
    \hspace{0.02\textwidth}
    \begin{minipage}[t]{0.30\textwidth}
        \centering
        \includegraphics[width=\linewidth]{20.5.png}
    \end{minipage}
}

\caption{Effect of occupancy-load support enforcement
\textbf{(a)} Removed predicted load versus horizon;
\textbf{(b)} raw versus support-consistent displacement error;
\textbf{(c)} raw versus support-consistent acceleration error;
\textbf{(d)} raw versus support-consistent peak-displacement error;
\textbf{(e)} distribution of absolute peak midspan acceleration.}
\label{fig:support_consistency_effect}
\end{figure*}


\subsection{Robustness to Structural Damping and Integration Time Step}
The 111-window sensitivity subset includes three windows from each held-out scenario. Across bridge damping ratios of 0\%, 1\%, 2\%, and 5\%, the ordering by midspan displacement and acceleration RMSE remains Support-Consistent Residual, Raw Residual, and Load-Matched Raw. Support-Consistent displacement RMSE changes from 4.009 to 3.913, 3.846, and 3.727~mm, respectively. At integration steps of 0.02, 0.01, and 0.005~s, its displacement RMSE is 4.0087, 4.0090, and 4.0081~mm, while acceleration RMSE remains near 0.0389~m/s$^2$. Hence, the ordering is robust to these damping settings and the baseline structural time step is numerically converged for this subset.



\subsection{Scope of Applicability and Future Research Directions}
\label{sec:scope_applicability}

The parameterized framework can potentially be adapted to highway bridges, overpasses, and viaducts by specifying span length, mass, modal frequencies and shapes, damping, and numerical integration settings. Its normalized traffic-to-bridge mapping permits spatiotemporal load fields to be transferred to an idealized structural span, while other traffic-loaded structures would require suitable modifications to the loading and structural models. Cross-country road-infrastructure monitoring research emphasizes geographically diverse data and region-specific model adaptation \cite{arya2024crosscountry}. Correspondingly, international applications would require locally representative vehicle classes, traffic behaviour and loading, together with bridge-specific structural properties.

The present study is limited to synthetic traffic and an idealized 318~m, five-mode reference bridge. Future investigations should consider other geometries and boundary conditions, detailed vehicle dynamics, road roughness, explicit future arrivals, and traffic-sampling resolution. Although damping and structural integration time-step sensitivities are evaluated here, deployment across structures and countries requires independent field traffic and bridge-response measurements, region-specific calibration, and numerical or experimental validation.

\FloatBarrier

\section{Conclusion}

This study developed a physics-guided framework that combines IDM/MOBIL traffic evolution, learned spatiotemporal residual correction, occupancy--load support enforcement, and multirate transfer to a reduced-order VBI model. Evaluation over 600 windows from 37 unseen synthetic traffic scenarios revealed that traffic-state accuracy and structural-response accuracy do not necessarily follow the same ranking: Residual Correction achieved the lowest aggregate-load error and improved occupied-cell estimation, whereas Physics-Based forecasting provided the lowest overall midspan displacement and acceleration errors.

Enforcing occupancy-load consistency reduced the maximum predicted midspan displacement from 191.12 to 40.61~mm and acceleration from 5.281 to 0.605~m/s$^2$. Against a frame-wise load-matched Raw Residual control, the Support-Consistent Residual reduced mean midspan displacement RMSE from 4.395 to 3.727~mm and acceleration RMSE from 0.0439 to 0.0321~m/s$^2$, with positive scenario-cluster confidence intervals. Thus, aggregate predicted load alone is insufficient to characterize downstream response; its spatial and lane-wise support also matters. The investigation remains a controlled numerical benchmark with an idealized structural model and a model-consistent traffic generator. Future work should examine traffic-model mismatch, future arrivals, more detailed bridge and vehicle dynamics, road roughness, traffic-resolution sensitivity, and independent field validation.

\section*{Acknowledgements}

This work was supported by the I-DAPT HUB FOUNDATION,
IIT (BHU), Varanasi, under the National Mission on
Interdisciplinary Cyber-Physical Systems (NM-ICPS),
Government of India.

\section*{Declaration of Competing Interest}

The authors declare that they have no known competing
financial interests or personal relationships that could
have appeared to influence the work reported in this paper.

\section*{Data Availability}

Data will be made available on request.

\FloatBarrier



\bibliographystyle{IEEEtran}
\bibliography{Ref}


\end{document}